\section{検証と残る変形問題}

全重みのhullについては、Poissonスキームの全局所環、形式弧の再中心化、
Hochschild一次差分と軌道切片を用いた。有限個の階数からの外挿ではない。
端点の全量子化はOre構造によって各ファイバーでAS正則かつCYとなる。
関数斜体の分類は先頭単項式の不変量と整数指数の可逆変換による。
PI次数の有限個の円分計算は、全位数の証明の確認として位置付ける。

VKの全次数の対称性保持族は、原始射に支持を持つ積補正の合成が零となることと、
平方零横断方向に関する有限な線形方程式によって構成する。
分母が原点で単元であることを確認し、得られた有理式について
全二千五百五十六個の結合則・導分方程式を多項式恒等式として照合した。
これにより五次元部分の二重構造を全次数で決定し、
三十五変数基底の中のイデアルも $(z_1,\ldots,z_{29},\gamma^2)$ と同定する。
純粋部分への還元では、商cochain複体の一次・二次コホモロジーの消滅を使って、
全次数でフィルトレーションを復元する。
補完接線は双対数上の具体的なOre代数から求め、全六項の合成を検査した。
同じ二次障害を、既存二次写像への代入と、実際のcochainからの独立した持上げ計算で求める。
通常の検証実行は入力と証明書を変更せず、完全な結果を凍結した証明書と比較する。

この結果は、標準六項補完像と対称性保持部分の位置関係を含む。
補完像は二重構造の横断長さ二の切片であり、滑らかな五変数部分の全てが
標準六項列へ補完できるという意味ではない。
別の例外列を許す半直交補完可能性の分類は、なお未決である。
導分付き障害複体を用いることで、外側の結合性障害線から内側の対称性障害への自然な比較写像を構成した。
一般の $HH^3$ 全体への写像は主張しない。
全重みで純粋部分は内在的に持ち上がり、同定を固定した相対変形は滑らかとなる。
純粋部分の $HH^1$ は一次元で、その生成元が全共鳴へ持ち上がるため、
接続写像は零であり、通常hullは純粋部分のhullと滑らかな因子に分解する。
一般の純粋hull $R_{N_m}$ の決定は未決である。
一般の重み・共鳴で、純粋部分の対称性保持方向を原始射の接合が何次まで切り詰めるかは、
この全次数の決定を拡張する次の問題となる。

端点では関数斜体を全パラメータで分類できたが、同じ斜体を持つ混合項 $H$ の方向を
射影圏や導来圏の不変量で区別する問題が残る。
公表稿XI\cite{VKPrimitive36}では、$H$ に依存しない中心開部分に対し、
境界を含む別の中心開部分を一般化Weyl代数として表示する。
その全中心は、$H\bmod x$ の巡回ノルム $P_H$ を用いた
$\Xi\Omega=P_H$ というアフィン改変であり、
標識付き境界方程式はこの剰余類を巡回作用まで復元する。
全中心ファイバーの単純加群と根基も切断間隔から決まり、
同じ関数斜体と同じ中心開部分を持ちながら異なる境界を持つ具体例を得る。

同稿では、偶数位数の特異部分上で根基を含む代数そのものを降下し、
二つの単純加群を分ける連結な二重被覆を決定する。
滑らかな中心部分のBrauer類の周期と一般点の指数は、共に冪根の位数 $\ell$ である。
従って完成局所的な行列環や非可換クレパント解消の表示は、
元の中心上での自己準同型環表示を意味しない。

原稿XIでは、$H=f(Q,u)+xG$ の全中心も決定し、
重み四・位数五で標識を忘れた全中心が非同型になる二つの混合項を示す。
この例では全アフィン代数のMorita同値とdg導来Morita同値も排除できる。
混合項を消せる部分族では全中心ファイバーを分類し、$f$ が $Q$ の倍数なら
最深点のファイバーと完備局所指数がそのまま保たれることも示す。
任意の混合項の深い境界と、混合項を動かした射影圏の同値分類は未決である。
CGG族内の導来分類を全miniversal族の分類へ拡大して述べることはしない。
